\chapter{Technologies and Tools}
\label{chap:technologies}

\section{Overview}

The project is implemented as a Python-based machine learning pipeline. The
current version uses PyTorch for supervised model training and PyKEEN for
knowledge graph embedding models.

\section{Programming Language and Environment}

\begin{table}[H]
\centering
\caption{Development environment}
\label{tab:development_environment}
\begin{tabularx}{\textwidth}{lX}
\toprule
\textbf{Item} & \textbf{Description} \\
\midrule
Language & Python 3.10 or newer \\
Operating system & Windows or Linux development environment \\
Deep learning framework & PyTorch with the CUDA 12.1 (\texttt{cu121}) wheel build \\
KGE framework & PyKEEN \\
Experiment output & CSV metrics, prediction files, saved model checkpoints \\
\bottomrule
\end{tabularx}
\end{table}

\section{PyTorch and CUDA Configuration}

The current implementation uses PyTorch as the deep learning framework for
the NFM classifier and tensor-based model training. GPU-enabled environments
install the official PyTorch wheel build compiled for NVIDIA CUDA 12.1,
identified by the \texttt{cu121} build tag. In a Jupyter or Colab notebook,
the installation command is:

\begin{lstlisting}[caption={Installing the PyTorch CUDA 12.1 wheel build}]
!pip install torch torchvision torchaudio --index-url https://download.pytorch.org/whl/cu121
\end{lstlisting}

For a standard terminal, the leading exclamation mark is omitted. The
equivalent command is:

\begin{lstlisting}[caption={Terminal installation command for PyTorch CUDA 12.1}]
pip install torch torchvision torchaudio --index-url https://download.pytorch.org/whl/cu121
\end{lstlisting}

The \texttt{cu121} index specifies the CUDA 12.1 wheel distribution. Because
the command does not pin numeric package releases, the exact versions of
\texttt{torch}, \texttt{torchvision}, and \texttt{torchaudio} are resolved
from the mutually compatible packages available on that index at installation
time.

\section{Main Libraries}

\begin{itemize}
    \item \textbf{PyTorch:} neural network implementation, tensor operations,
    model training, and checkpoint saving.

    \item \textbf{PyKEEN:} training and inference for knowledge graph embedding
    models such as DistMult and CompGCN.

    \item \textbf{NumPy:} numerical operations and array processing.

    \item \textbf{pandas:} tabular data loading, merging, and output writing.

    \item \textbf{scikit-learn:} preprocessing utilities and evaluation support.

    \item \textbf{tqdm:} progress visualization during experiment runs.
\end{itemize}

\section{Project Structure}

\begin{table}[H]
\centering
\caption{Important source files}
\label{tab:source_files}
\begin{tabularx}{\textwidth}{lX}
\toprule
\textbf{Path} & \textbf{Purpose} \\
\midrule
\texttt{src/kge\_nfm.py} & Main PyTorch experiment runner \\
\texttt{src/data\_utils.py} & Dataset loading, fold loading, KG loading, and descriptor loading \\
\texttt{src/kge.py} & KGE training, scoring, and embedding extraction utilities \\
\texttt{src/nfm.py} & NFM model definition and training loop \\
\texttt{src/metrics\_utils.py} & ROC-AUC and PR-AUC helper functions \\
\texttt{src/utils.py} & Seed setting, device selection, and output directory helpers \\
\texttt{baseline/} & Baseline models and exploratory experiments \\
\texttt{output/} & Experiment results, model files, predictions, and metric curves \\
\bottomrule
\end{tabularx}
\end{table}


\section{Example Command}

\begin{lstlisting}[caption={Example KGE-NFM experiment command}]
python src/kge_nfm.py \
  --dataset yamanishi_08 \
  --split warm_start_1_10 \
  --folds 10 \
  --device auto \
  --embedding-dim 400 \
  --kge-epochs 50 \
  --nfm-epochs 2000 \
  --output-dir output/kge_nfm_torch
\end{lstlisting}
